\documentclass{beamer}
\usepackage[utf8]{inputenc}
\usepackage{lmodern}
\usepackage[ngerman]{babel}
\usepackage{amsmath,amsfonts,amssymb}
\usepackage{graphicx}
\usetheme{Copenhagen}
\useinnertheme{rectangles}
\beamertemplatenavigationsymbolsempty
\setbeamertemplate{footline}[frame number]
\title{YUV.KA \\ Praxis der Softwareentwicklung 2011/2012}
\author{Max Wagner $\cdot$ Patrick Gemander $\cdot$ Sebastian Ullrich \\ Michael Vollmer $\cdot$ Robert Hangu $\cdot$ Daniel Lebert}
\titlegraphic{\includegraphics[height=2.5cm]{logo.png}}
\begin{document}
\frame{\titlepage}

\frame{\tableofcontents}

\section{Einleitung}
\frame{\frametitle{Aufgabenstellung}
Multimedia-Framework zur Evaluation von Videoencodern \pause
\begin{itemize}
	\item Einlesen von unkomprimierten Videos \pause
	\item Manipulation durch Syntheseoptionen \pause
	\item Enkodierung der Videodateien \pause
	\item Vergleich mit Referenzvideos
\end{itemize}
}

\section{Benutzeroberfläche}
\subsection{Hauptanzeige}
\frame{
\frametitle{Benutzeroberfläche I}
\begin{columns}
        \column{.55\textwidth}
                \fbox{\pgfimage[width=\textwidth]{main-screen}} \pause
        \column{.45\textwidth}
                \begin{itemize}
		\item Eingabe von Videos, Bildern, Farbe, Noise \pause
		\item Speichern von Videos und Pipeline \pause
		\item Anzeigen von Videos \pause
		\item Abspielgeschwindigkeit verändern \pause
		\item Im Video springen \pause
		\item Videos zusammenführen \pause
		\item Videos trennen \pause
		\item Videos weichzeichnen
                \end{itemize}
\end{columns}
}

\frame{
\frametitle{Benutzeroberfläche II}
\begin{columns}
        \column{.55\textwidth}
                \fbox{\pgfimage[width=\textwidth]{main-screen}}
        \column{.45\textwidth}
                \begin{itemize}
		\item Videos invertieren \pause
		\item Kontrast, Sättigung, Helligkeit verändern \pause
		\item Verzögern eines Videos \pause
		\item Differenzvideos bilden \pause
		\item Videos per Testencoder enkodieren \pause
		\item Analyse der Videos \pause
		\item Video mit Overlayoptionen abspielen
                \end{itemize}
\end{columns}
}

\subsection{Analyseansicht}
\frame{
\frametitle{Analyseknoten I}
\begin{columns}
        \column{.55\textwidth}
                \fbox{\pgfimage[width=\textwidth]{diagram-screen}} \pause
        \column{.45\textwidth}
                \begin{itemize}
		\item Inter/Intra-Block-Verhältnis \pause
		\item Helligkeitverlauf \pause
		\item Histogramm \pause
		\item PSNR
                \end{itemize}
\end{columns}
}

\frame{
\frametitle{Analyseknoten II}
\begin{columns}
        \column{.55\textwidth}
                \fbox{\pgfimage[width=\textwidth]{diagram-screen}}
        \column{.45\textwidth}
                \begin{itemize}
		\item Differenzen der Pixelfarben \pause
		\item Differenz der Encoder-Entscheidungen \pause
		\item Anzahl der Artefakte
                \end{itemize}
\end{columns}
}

\section{Wunschkriterien}
\frame{
\frametitle{Wunschkriterien I}
\begin{columns}
        \column{.55\textwidth}
                \fbox{\pgfimage[width=\textwidth]{main-screen}} \pause
        \column{.45\textwidth}
                \begin{itemize}
		\item Plugins \pause
		\item Quantisieren der Farbpalette \pause
		\item Undo/Redo-Funktion \pause
		\item Tabs im Interface \pause
		\item asynchrone, nichtblockende UI
                \end{itemize}
\end{columns}
}

\frame{
\frametitle{Wunschkriterien II} \pause
                \begin{itemize}
		\item autom. Auflösungserkennung \pause
		\item Multithreading \pause
		\item Motionvektoren verarbeiten
                \end{itemize}
}

\section{Systemmodell}
\frame{
\frametitle{Systemmodell}
\begin{columns}
        \column{.40\textwidth}
                \fbox{\pgfimage[width=\textwidth]{Layers}} \pause
        \column{.60\textwidth}
                \begin{description}
		\item[View Model] eingelesenes Video mit Log-Daten \pause
		\item[Encoder] der zu testende Encoder \pause
		\item[Pipeline] Aufbau und Abarbeitung der Pipeline, bestend aus \pause
			\begin{itemize}
				\item Knoten-Klassen \pause
				\item Pipeline Driver \pause
			\end{itemize}
		\item[ViewModel] Abarbeitung von Benutzereingaben \pause
		\item[View] Aufbau und Anzeige der Benutzeroberfläche
                \end{description}
\end{columns}
}

\section{Fragen}
\frame{\frametitle{Fragen}}

\end{document}